\subsection{\texorpdfstring{THE ISOMORPHISMS \(\operatorname{Hom}_{\mathbf{B}}(\mathbf{E} \otimes_{\mathbf{A}} \mathbf{F}, \mathbf{G}) \to \operatorname{Hom}_{\mathbf{A}}(\mathbf{F}, \operatorname{Hom}_{\mathbf{B}}(\mathbf{E}, \mathbf{G}))\) AND \(\operatorname{Hom}_{\mathbf{C}}(\mathbf{E} \otimes_{\mathbf{A}} \mathbf{F}, \mathbf{G}) \to \operatorname{Hom}_{\mathbf{A}}(\mathbf{E}, \operatorname{Hom}_{\mathbf{C}}(\mathbf{F}, \mathbf{G}))\)}{THE ISOMORPHISMS \textbackslash operatorname\{Hom\}\_\{\textbackslash mathbf\{B\}\}(\textbackslash mathbf\{E\} \textbackslash otimes\_\{\textbackslash mathbf\{A\}\} \textbackslash mathbf\{F\}, \textbackslash mathbf\{G\}) \textbackslash to \textbackslash operatorname\{Hom\}\_\{\textbackslash mathbf\{A\}\}(\textbackslash mathbf\{F\}, \textbackslash operatorname\{Hom\}\_\{\textbackslash mathbf\{B\}\}(\textbackslash mathbf\{E\}, \textbackslash mathbf\{G\})) AND \textbackslash operatorname\{Hom\}\_\{\textbackslash mathbf\{C\}\}(\textbackslash mathbf\{E\} \textbackslash otimes\_\{\textbackslash mathbf\{A\}\} \textbackslash mathbf\{F\}, \textbackslash mathbf\{G\}) \textbackslash to \textbackslash operatorname\{Hom\}\_\{\textbackslash mathbf\{A\}\}(\textbackslash mathbf\{E\}, \textbackslash operatorname\{Hom\}\_\{\textbackslash mathbf\{C\}\}(\textbackslash mathbf\{F\}, \textbackslash mathbf\{G\}))}}

Let E be a right A-module, F a left A-module, G a Z-module and H the Z-module of mappings \(f: E \times F \rightarrow G\) which are Z-bilinear and satisfy

\[
f (x \lambda , y) = f (x, \lambda y) \quad \text { for } x \in \mathrm{E}, y \in \mathrm{F}, \lambda \in \mathrm{A}.\tag{1}
\]

It has been seen (§ 3, no. 1, Proposition 1) that there exists a canonical Z-module homomorphism

\[
\mathrm{H} \rightarrow \operatorname{Hom} _ {\mathbf {Z}} (\mathrm{E} \otimes_ {\mathrm{A}} \mathrm{F}, \mathrm{G}).\tag{2}
\]

On the other hand, a left A-module structure has been defined on \(\mathrm{Hom}_{\mathbf{Z}}(\mathrm{E},\mathrm{G})\) and a right A-module structure on \(\mathrm{Hom}_{\mathbf{Z}}(\mathrm{F},\mathrm{G})\) (§ 3, no. 3); we may therefore consider the Z-modules \(\mathrm{Hom}_{\mathbf{A}}(\mathrm{E},\mathrm{Hom}_{\mathbf{Z}}(\mathrm{F},\mathrm{G}))\) and \(\mathrm{Hom}_{\mathbf{A}}(\mathrm{F},\mathrm{Hom}_{\mathbf{Z}}(\mathrm{E},\mathrm{G}))\) . A mapping \(f\) of \(\mathbf{E} \times \mathbf{F}\) into \(\mathbf{G}\) is canonically identified with a mapping of \(\mathbf{E}\) into the set \(\mathbf{G}^{\mathbf{F}}\) of mappings of \(\mathbf{F}\) into \(\mathbf{G}\) (Set Theory, II, § 5, no. 2); by expressing the fact that the latter mapping belongs to \(\mathrm{Hom}_{\mathbf{A}}(\mathrm{E},\mathrm{Hom}_{\mathbf{Z}}(\mathrm{F},\mathrm{G}))\) , we obtain precisely the fact that \(f\) is biadditive and conditions (1); whence there is a canonical isomorphism

\[
\mathrm{H} \rightarrow \operatorname{Hom} _ {\mathbf {A}} (\mathrm{E}, \operatorname{Hom} _ {\mathbf {Z}} (\mathrm{F}, \mathrm{G}))\tag{3}
\]

and similarly there is defined a canonical isomorphism

\[
\mathrm{H} \rightarrow \operatorname{Hom} _ {\mathbf {A}} (\mathrm{F}, \operatorname{Hom} _ {\mathbf {Z}} (\mathrm{E}, \mathrm{G})).\tag{4}
\]

Suppose now that E and G also have left (resp. right) B-module structures and that the A-module and B-module structures on E are compatible. Then \(E \otimes_{A} F\) has canonically a left (resp. right) B-module structure ( \(\S 3\) , no. 4) and on the other hand \(\operatorname{Hom}_{\mathbf{B}}(\mathbf{E}, \mathbf{G})\) has canonically a left A-module structure ( \(\S 1\) , no. 14). We may therefore consider the Z-modules \(\operatorname{Hom}_{\mathbf{B}}(\mathbf{E} \otimes_{\mathbf{A}} \mathbf{F}, \mathbf{G})\) and \(\operatorname{Hom}_{\mathbf{A}}(\mathbf{F}, \operatorname{Hom}_{\mathbf{B}}(\mathbf{E}, \mathbf{G}))\) , which are submodules of \(\operatorname{Hom}_{\mathbf{Z}}(\mathbf{E} \otimes_{\mathbf{A}} \mathbf{F}, \mathbf{G})\) and \(\operatorname{Hom}_{\mathbf{A}}(\mathbf{F}, \operatorname{Hom}_{\mathbf{Z}}(\mathbf{E}, \mathbf{G}))\) respectively ( \(\S 2\) , no. 1, Theorem 2). We examine under what condition a mapping \(f \in H\) has as image under the isomorphisms (2) and (4) an element of \(\operatorname{Hom}_{\mathbf{B}}(\mathbf{E} \otimes_{\mathbf{A}} \mathbf{F}, \mathbf{G})\) and an element of \(\operatorname{Hom}_{\mathbf{A}}(\mathbf{E}, \operatorname{Hom}_{\mathbf{B}}(\mathbf{F}, \mathbf{G}))\) respectively; in each of the two cases we find the same condition

\[
\begin{array}{r l} f (\beta x, y) & = \beta f (x, y) \\ (\text { resp.   } f (x \beta , y) & = f (x, y) \beta) \end{array}
\]

for \(x \in \mathbf{E}, y \in \mathbf{F}, \beta \in \mathbf{B}\) .

Similarly, suppose that F and G are left (resp. right) C-modules and that the A-module and C-module structures on F are compatible. Then, for a mapping \(f \in H\) to have as image under (2) or (3) an element of \(\operatorname{Hom}_{\mathbf{C}}(\mathbf{E} \otimes_{\mathbf{A}} \mathbf{F}, \mathbf{G})\) or \(\operatorname{Hom}_{\mathbf{A}}(\mathbf{E}, \operatorname{Hom}_{\mathbf{C}}(\mathbf{F}, \mathbf{G}))\) respectively, it is necessary and sufficient that it satisfy the same condition

\[
\begin{array}{c} f (x, \gamma y) = \gamma f (x, y) \\ (\text { resp. } f (x, y \gamma) = f (x, y) \gamma) \end{array}
\]

for \(x \in \mathbf{E}, y \in \mathbf{F}, \gamma \in \mathbf{C}\) .

We have therefore established the following result (in the notation introduced above):

PROPOSITION 1. (a) Let E be a (B, A)-bimodule, F a left A-module and G a left B-module. For every mapping \(g \in \operatorname{Hom}_{\mathbf{B}}(\mathbf{E} \otimes_{\mathbf{A}} \mathbf{F}, \mathbf{G})\) , let \(g'\) be the mapping of F into \(\operatorname{Hom}_{\mathbf{B}}(\mathbf{E}, \mathbf{G})\) defined by \((g'(y))(x) = g(x \otimes y)\) for \(x \in \mathbf{E}, y \in \mathbf{F}\) . The mapping \(g \mapsto g'\) is an isomorphism

\[
\beta : \operatorname{Hom} _ {\mathbf {B}} (\mathbf {E} \otimes_ {\mathbf {A}} \mathbf {F}, \mathbf {G}) \rightarrow \operatorname{Hom} _ {\mathbf {A}} (\mathbf {F}, \operatorname{Hom} _ {\mathbf {B}} (\mathbf {E}, \mathbf {G})).\tag{5}
\]

\begin{enumerate}
\def\labelenumi{(\alph{enumi})}
\setcounter{enumi}{1}
\tightlist
\item
  Let \(\mathbf{E}\) be a right \(\mathbf{A}\) -module, \(\mathbf{F}\) an \((\mathbf{A}, \mathbf{C})\) -bimodule and \(\mathbf{G}\) a right \(\mathbf{C}\) -module. For every mapping \(h \in \operatorname{Hom}_{\mathbf{C}}(\mathbf{E} \otimes_{\mathbf{A}} \mathbf{F}, \mathbf{G})\) , let \(h'\) be the mapping of \(\mathbf{E}\) into \(\operatorname{Hom}_{\mathbf{C}}(\mathbf{F}, \mathbf{G})\) defined by \((h'(x))(y) = h(x \otimes y)\) for \(x \in \mathbf{E}\) , \(y \in \mathbf{F}\) . The mapping \(h \mapsto h'\) is an isomorphism
\end{enumerate}

\[
\gamma : \operatorname{Hom} _ {\mathbf {C}} (\mathrm{E} \otimes_ {\mathbf {A}} \mathrm{F}, \mathrm{G}) \rightarrow \operatorname{Hom} _ {\mathbf {A}} (\mathrm{E}, \operatorname{Hom} _ {\mathbf {C}} (\mathrm{F}, \mathrm{G})).\tag{6}
\]

In particular B and C may be taken to be a subring \(\Gamma\) of the centre of the ring A; then for every \(\Gamma\) -module G, the three \(\Gamma\) -modules

\[
\operatorname{Hom} _ {\Gamma} (\mathrm{E} \otimes_ {\mathrm{A}} \mathrm{F}, \mathrm{G}), \quad \operatorname{Hom} _ {\mathrm{A}} (\mathrm{E}, \operatorname{Hom} _ {\Gamma} (\mathrm{F}, \mathrm{G})), \quad \operatorname{Hom} _ {\mathrm{A}} (\mathrm{F}, \operatorname{Hom} _ {\Gamma} (\mathrm{E}, \mathrm{G}))
\]

are canonically isomorphic to the \(\Gamma\) -module of \(\Gamma\) -bilinear mappings of \(\mathbf{E} \times \mathbf{F}\) into \(\mathbf{G}\) which satisfy (1). More particularly:

COROLLARY. If C is a commutative ring and E, F, G three C-modules, then the C-modules

\[
\begin{array}{l l} \operatorname{Hom} _ {\mathbf {c}} (\mathrm{E} \otimes_ {\mathbf {c}} \mathrm{F}, \mathrm{G}), & \operatorname{Hom} _ {\mathbf {c}} (\mathrm{E}, \operatorname{Hom} _ {\mathbf {c}} (\mathrm{F}, \mathrm{G})), \\ \operatorname{Hom} _ {\mathbf {c}} (\mathrm{F}, \operatorname{Hom} _ {\mathbf {c}} (\mathrm{E}, \mathrm{G})), & \mathcal {L} _ {2} (\mathrm{E}, \mathrm{F}; \mathrm{G}) \end{array}
\]

are canonically isomorphic.
% semantic-fragments: legacy-input-seam
\relax{}
